
\begin{frame}
\frametitle{Bayesian Calibration using Universal Scaled Spectra}

\begin{columns}[T]

%====================================================
% LEFT
%====================================================

\column{0.48\textwidth}

\centering

\Large

\[
\boxed{
\text{Universal Scaled Spectra}
}
\]

\vspace{0.25cm}

\[
\Downarrow
\]

\vspace{0.15cm}

\[
\boxed{
\text{Hydrodynamic Model}
}
\]

\vspace{0.25cm}

\[
\Downarrow
\]

\vspace{0.15cm}

\[
\boxed{
\text{Gaussian Process Emulator}
}
\]

\vspace{0.25cm}

\[
\Downarrow
\]

\vspace{0.15cm}

\[
\boxed{
\text{Bayesian Inference}
}
\]

%====================================================
% RIGHT
%====================================================

\column{0.52\textwidth}

\small

\textbf{Calibration procedure}

\begin{itemize}
    \item Universal scaled spectra are used as experimental observables.
    \item Event-by-event hydrodynamic simulations are replaced by Gaussian Process emulators.
    \item Bayesian inference constrains the QGP model parameters.
\end{itemize}

\vspace{0.25cm}

\textbf{This work}

\begin{itemize}
    \item Keep the original Bayesian calibration framework.
    \item Replace the diagonal experimental covariance matrix by correlated covariance models (CS and RBF).
    \item Quantify the impact on the inferred QGP parameters.
\end{itemize}

\end{columns}

\end{frame}